% Generated by scripts/build-paper.py — edit paper/src/survey.src.md
\section{Introduction}\label{sec:introduction}

Many software decisions are small and bounded: route a ticket, accept or escalate a judgment, pick the next tool, code a variable from a narrative. Large language models can make such decisions, but they return strings that must be parsed, validated and retried, and their stated confidence is usually a verbalised estimate rather than a probability the program can threshold \citep{arxiv2207_05221,arxiv2503_15850}. Typed decision models change the contract. The caller declares the admissible answers; the model returns one of them together with a probability for every option; no free text is generated, so a parse failure is impossible by construction.

TypeSafe AI's Jev, released in early access on 15 September 2026, made this contract prominent \citep{typesafe2026launch}. Its launch materials describe a “new class of frontier models” optimised with “Reinforcement Learning for Calibrated Decisions” and report end-to-end responses of 70–500 ms at \$0.042 per million input tokens \citep{typesafe2026launch,typesafe2026models}. Between 19 and 23 September 2026, 17 arXiv preprints evaluated Jev, built open Jev-like models, or used typed decisions inside larger systems. The claims in this literature range from schema validity to calibration, cost, latency and end-task success; the measurement protocols range from 20 curated scientific choices to 499,500 crash narratives.

A public survey draft, \emph{Decisions, Not Tokens}, dated 21 September 2026, already places typed decisions in the history of classification, calibration, selective prediction and routing, and proposes a five-dimensional taxonomy and four evaluation layers \citep{decisionsnottokens2026}. Its bibliography at the pinned commit contains none of the core preprints reviewed here. This survey therefore does something different: it is an \textbf{empirical survey and evidence audit} of the first wave, organised so that incomparable numbers stay apart.

We keep three objects distinct throughout: (A) the \textbf{commercial model}, TypeSafe's hosted Jev, whose architecture and training data are undisclosed; (B) the \textbf{paradigm}, Jev-like typed decision models and readouts built by others; and (C) \textbf{systems} whose results depend on typed decisions. A result about one is not evidence about another.

We ask six questions. \textbf{RQ1 (definition and lineage):} how do runtime-defined finite-option decisions relate to classifiers, zero-shot labelling, rerankers, reward models and constrained generation? \textbf{RQ2 (probability quality):} what do returned probabilities and confidence mean, and where do they hold? \textbf{RQ3 (efficiency attribution):} how much of a speed or cost gain comes from the readout, and how much from model size, batching, prefix sharing, caching or the network? \textbf{RQ4 (selective control):} which decisions can be delegated, and when should software escalate, abstain or defer to code? \textbf{RQ5 (robustness):} how do option names, order, rubric binding, missing options, language and distribution shift change decisions? \textbf{RQ6 (openness):} what do alternatives release, and how much of the literature survives a common protocol?

Our contributions are concrete artefacts rather than claims of priority:

\begin{enumerate}[leftmargin=*,itemsep=2pt]
  \item A versioned evidence base of 71 references and 75 evidence records. Each record carries a source locator, evidence type, test level, measurement scope, sample size and model version, with explicit reasons where a value is missing (§3).
  \item A synthesis into seven findings, each listing supporting \emph{and} qualifying evidence (§7). Study families are synthesised once, and nothing is pooled across incompatible scopes.
  \item An implementation lineage of six readout families and an openness audit that separates code, weights, data, raw predictions, recomputability and reproduction (§6, §9).
  \item A catalogue of failure modes with sources and mitigations (§8), and a proposed evaluation protocol whose first experiment targets the intersection of name binding, calibration and escalation (§10).
  \item A public website, generated README, BibTeX and CSV exports, all produced from the same data by scripts with validation and link checks.
\end{enumerate}
